\section{Display: block, inline, dan inline-block}

Properti \code{display} menentukan bagaimana elemen berpartisipasi dalam \textbf{flow layout}---model tata letak default browser di mana elemen disusun secara berurutan dari atas ke bawah dan kiri ke kanan. Setiap elemen HTML memiliki nilai display default: beberapa bersifat \textbf{block} (memenuhi satu baris penuh), yang lain bersifat \textbf{inline} (mengalir di dalam baris teks). Properti \code{display} memungkinkan pengembang mengubah perilaku default ini sesuai kebutuhan desain \cite{mdnCss}.

\begin{table}[htbp]
\centering
\begin{tabular}{|P{2.5cm}|P{3cm}|P{2.5cm}|P{4cm}|}
\hline
\textbf{Nilai} & \textbf{Perilaku} & \textbf{W/H} & \textbf{Contoh Elemen Default} \\
\hline
\code{block} & Satu baris penuh; mulai baris baru & Diterima & \code{<div>}, \code{<p>}, \code{<h1>}, \code{<section>} \\
\hline
\code{inline} & Mengalir dalam teks; sebaris & Diabaikan & \code{<span>}, \code{<a>}, \code{<strong>}, \code{<em>} \\
\hline
\code{inline-block} & Sebaris, dapat diberi ukuran & Diterima & Tidak ada default; harus diset \\
\hline
\code{none} & Elemen hilang dari layout & --- & Disembunyikan sepenuhnya \\
\hline
\end{tabular}
\caption{Perbandingan nilai display: block, inline, inline-block, dan none}
\end{table}

Elemen \textbf{block} secara default memakai seluruh lebar kontainer induknya dan memulai baris baru. Properti \code{width}, \code{height}, \code{margin}, dan \code{padding} berlaku sepenuhnya pada elemen block. Elemen \textbf{inline} hanya memakai ruang sebesar kontennya dan mengalir di dalam baris teks bersama konten lain. Pada elemen inline, \code{width} dan \code{height} \textbf{tidak berlaku}, dan margin/padding vertikal tidak mempengaruhi layout elemen lain. Elemen \textbf{inline-block} menggabungkan keduanya: mengalir di dalam baris seperti inline, tetapi menerima \code{width}, \code{height}, dan margin/padding penuh seperti block \cite{webdevCss}.

Properti \code{display: none;} menghapus elemen dari layout sepenuhnya---elemen tidak terlihat dan tidak mengambil ruang. Berbeda dengan \code{visibility: hidden;} yang menyembunyikan elemen secara visual tetapi tetap mempertahankan ruangnya di layout. Pilihan antara keduanya tergantung pada apakah ruang kosong harus tetap ada atau tidak. Untuk aksesibilitas, \code{display: none} juga membuat elemen tidak terbaca oleh pembaca layar, sedangkan \code{visibility: hidden} masih dapat terbaca tergantung implementasi \cite{mdnCss}.

